\documentclass[11pt]{article}
\usepackage[margin=1in]{geometry}
\usepackage{amsmath,amssymb,amsthm}
\usepackage{graphicx}
\usepackage{booktabs}
\usepackage[hidelinks]{hyperref}

\newtheorem{theorem}{Theorem}
\newtheorem{lemma}{Lemma}
\newcommand{\T}{\mathbb T}
\newcommand{\N}{\mathbb N}

\title{A cosecant certificate for the bidisc weak-product constant}
\author{Automated mathematics run \texttt{fa\_banach\_001}, lane 0}
\date{21 July 2026}

\begin{document}
\maketitle

\begin{center}
\fbox{\parbox{0.91\textwidth}{\textbf{Status: candidate partial result,
likely valid.}  The result below improves a quantitative obstruction but does
not resolve the equality of the two spaces.}}
\end{center}

\section{Source question}

Bampouras and Brevig ask whether, for fixed $d\geq2$, there is a constant
$C_d>0$ such that
\[
 \|f\|_{H^1(\T^d)}\geq C_d\|f\|_{W(\T^d)},
 \qquad W(\T^d)=H^2(\T^d)\odot H^2(\T^d),
\]
for every $f\in W(\T^d)$; equivalently, whether $H^1(\T^d)$ and
$W(\T^d)$ are equal as sets \cite{BB}.  The question occurs on page 3 of the
source paper.

\begin{figure}[ht]
 \centering
 \includegraphics[width=\textwidth]{figures/open_problem_crop.png}
 \caption{The open question in arXiv:2310.20346, page 3.}
\end{figure}

For the bidisc, put
\[
 K_2:=\sup_{0\neq f\in W(\T^2)}
 \frac{\|f\|_{W(\T^2)}}{\|f\|_{H^1(\T^2)}}\in[1,\infty].
\]
Thus the best possible constant in the displayed inequality is $K_2^{-1}$
when $K_2<\infty$.  The quadratic construction used previously gives a lower
bound about $1.16064$ for $K_2$ \cite{Brevig}; the source paper also discusses
the earlier Ortega--Cerd\`a--Seip obstruction \cite{OS}.

\section{Proof intuition}

The source paper turns the weak-product dual norm into the operator norm of a
small Hankel operator.  For a homogeneous bidisc symbol, that Hankel operator
is a direct sum of finite Hankel matrices.  We choose the symbol coefficients
from a discrete cosecant matrix which is exactly orthogonal.  Every required
Hankel block is then a submatrix of an isometry, giving a norm-one dual
certificate at once.

Pairing this certificate with a homogeneous Dirichlet polynomial yields a
quantity asymptotic to $(2/\pi)\log N$.  The $H^1$ norm of that polynomial is
the usual Dirichlet-kernel $L^1$ norm, asymptotic to
$(4/\pi^2)\log N$.  The quotient of the two leading constants is $\pi/2$.
This supplies a quantitative partial result only; it neither proves nor
disproves that $K_2$ is finite.

\section{Statement}

\begin{theorem}\label{thm:main}
Let $W(\T^2)=H^2(\T^2)\odot H^2(\T^2)$.  Then
\[
 K_2\geq \frac{\pi}{2}.
\]
Consequently, if a constant $C_2>0$ satisfies
$\|f\|_{H^1(\T^2)}\geq C_2\|f\|_{W(\T^2)}$ for every $f\in W(\T^2)$, then
\[
 C_2\leq \frac{2}{\pi}.
\]
More precisely, for $N\geq2$ define
\[
 f_N(z_1,z_2)=\sum_{k=0}^{N-1}z_1^kz_2^{N-1-k}
\]
and
\[
 S_N=\frac1N\sum_{k=0}^{N-1}
 \csc\!\left(\frac{\pi(k+\tfrac12)}{N}\right).
\]
Then $\|f_N\|_{W(\T^2)}\geq S_N$ and
\[
 \frac{S_N}{\|f_N\|_{H^1(\T^2)}}\longrightarrow\frac\pi2.
\]
\end{theorem}

\section{Proof}

We first record the finite orthogonality identity responsible for the
certificate.

\begin{lemma}\label{lem:orthogonal}
For $N\geq2$, the real $N\times N$ matrix
\[
 C_{r,s}=\frac{1}{N\sin\!\left(\pi(r+s+\tfrac12)/N\right)},
 \qquad 0\leq r,s\leq N-1,
\]
is orthogonal.
\end{lemma}

\begin{proof}
Use the standard identities, obtained by logarithmically differentiating the
finite sine product,
\begin{align*}
 \sum_{s=0}^{N-1}\cot\left(x+\frac{s\pi}{N}\right)
   &=N\cot(Nx),\\
 \sum_{s=0}^{N-1}\csc^2\left(x+\frac{s\pi}{N}\right)
   &=N^2\csc^2(Nx).
\end{align*}
For distinct rows $r\neq t$, set
$u_s=\pi(r+s+\tfrac12)/N$ and
$v_s=\pi(t+s+\tfrac12)/N$.  Since
\[
 \frac1{\sin u_s\sin v_s}
 =\frac{\cot u_s-\cot v_s}{\sin(v_s-u_s)},
\]
the inner product of rows $r$ and $t$ vanishes: both cotangent sums are
$N\cot(\pi(r+\tfrac12))=N\cot(\pi(t+\tfrac12))=0$.
For a row with itself, the second identity gives
\[
 \sum_{s=0}^{N-1}\csc^2 u_s
 =N^2\csc^2\!\left(\pi(r+\tfrac12)\right)=N^2.
\]
The prefactor $N^{-2}$ completes the proof.
\end{proof}

Now put $m=N-1$ and consider the $m$-homogeneous symbol
\[
 \psi_N(z_1,z_2)=\sum_{k=0}^{m}a_kz_1^kz_2^{m-k},
 \qquad
 a_k=\frac{1}{N\sin\!\left(\pi(k+\tfrac12)/N\right)}.
\]
The dual of $W(\T^2)$ is isometric to the bounded small Hankel operators
\cite{BB}.  For completeness, on the homogeneous input space of degree $n$,
$0\leq n\leq m$, the matrix of the Hankel form associated with $\psi_N$ is
\[
 A_n=(a_{i+j})_{\substack{0\leq i\leq m-n\\0\leq j\leq n}}.
\]
Indeed, the product of the basis monomials
$z_1^jz_2^{n-j}$ and $z_1^iz_2^{m-n-i}$ has $z_1$-degree $i+j$.
The matrix $A_n$ is exactly the submatrix of $C$ formed by rows
$0,\ldots,m-n$ and columns $0,\ldots,n$.  Lemma~\ref{lem:orthogonal} therefore
gives $\|A_n\|_{\mathrm{op}}\leq1$ for every $n$.  The full Hankel operator is
the orthogonal direct sum of these blocks (and is zero on input degrees larger
than $m$), so
\[
 \|\psi_N\|_{W(\T^2)^*}\leq1.
\]
Duality now yields
\[
 \|f_N\|_{W(\T^2)}
 \geq |\langle f_N,\psi_N\rangle|
 =\sum_{k=0}^{N-1}a_k=S_N.                                      \tag{1}
\]

It remains to compare two elementary logarithmic asymptotics.  Set
$x_k=(k+\tfrac12)/N$.  The function
\[
 h(x)=\csc(\pi x)-\frac1{\pi x}-\frac1{\pi(1-x)}
\]
extends continuously to $[0,1]$.  Consequently, midpoint summation and a
change of index give
\begin{align*}
 S_N
 &=\frac1N\sum_{k=0}^{N-1}
 \left(\frac1{\pi x_k}+\frac1{\pi(1-x_k)}+h(x_k)\right)\\
 &=\frac2\pi\sum_{k=0}^{N-1}\frac1{k+\tfrac12}+O(1)
 =\frac2\pi\log N+O(1).                                      \tag{2}
\end{align*}

On the torus,
\[
 |f_N(e^{i\theta_1},e^{i\theta_2})|
 =\left|\sum_{k=0}^{N-1}e^{ik(\theta_1-\theta_2)}\right|.
\]
Hence normalized Haar integration and the Dirichlet-kernel formula give
\[
 \|f_N\|_{H^1(\T^2)}
 =\frac1\pi\int_0^\pi\frac{|\sin(Nx)|}{\sin x}\,dx
 =\frac2\pi\int_0^{\pi/2}\frac{|\sin(Nx)|}{\sin x}\,dx.       \tag{3}
\]
On $(0,\pi/2]$, the difference $1/\sin x-1/x$ is bounded by a constant
times $x$.  Replacing $1/\sin x$ by $1/x$ in (3) therefore changes the
integral by $O(1)$.  After $u=Nx$,
\[
 \|f_N\|_{H^1(\T^2)}
 =\frac2\pi\int_0^{N\pi/2}\frac{|\sin u|}{u}\,du+O(1).
\]
The periodic function $|\sin u|-2/\pi$ has mean zero and a bounded
primitive.  Integration by parts on $[1,R]$ thus shows
\[
 \int_0^R\frac{|\sin u|}{u}\,du=\frac2\pi\log R+O(1).
\]
Therefore
\[
 \|f_N\|_{H^1(\T^2)}=\frac4{\pi^2}\log N+O(1).                 \tag{4}
\]
Combining (1)--(4) proves
\[
 \liminf_{N\to\infty}
 \frac{\|f_N\|_{W(\T^2)}}{\|f_N\|_{H^1(\T^2)}}
 \geq\frac{(2/\pi)}{(4/\pi^2)}=\frac\pi2,
\]
which proves Theorem~\ref{thm:main}.

\section{Verification notes}

The included script \texttt{code/verify\_cosecant\_certificate.py} checks, for
several $N$, that $CC^T$ is the identity to floating-point precision, every
block $A_n$ has operator norm at most one, and the numerical certificate ratio
approaches $\pi/2$.  These tests are sanity checks only; Lemma~\ref{lem:orthogonal}
and the asymptotic calculation are complete analytic proofs.

The most important human checks are:
\begin{enumerate}
 \item the convention for the small Hankel pairing and the indices of $A_n$;
 \item that all homogeneous blocks, including the endpoint degrees, are
 covered by the submatrix argument;
 \item the normalization of Haar measure in the Dirichlet-kernel asymptotic.
\end{enumerate}

\section{Limitations and novelty status}

The theorem only bounds a hypothetical embedding constant.  It does not show
that $K_2$ is finite, so the source problem remains open.  A digit-separated
lacunary product test did not amplify the finite gain: extra Hankel
interactions increased the dual norm faster than the observed $H^1$ gain.

A bounded search on 21 July 2026 inspected the source paper and Brevig's
earlier paper \cite{Brevig}.  It also checked arXiv:2101.00763 and exact
arXiv/web queries containing ``polydisc Nehari'', ``cosecant matrix'',
``Dirichlet kernel'', and ``$\pi/2$''.  No matching construction or bound was
found.  The $N=3$ member
of the family recovers the earlier quadratic certificate in \cite{Brevig},
which is evidence that the present construction is a natural asymptotic
extension, not evidence of definitive novelty.  Novelty confidence is
therefore moderate and a specialist literature review is recommended.

\begin{thebibliography}{9}
\bibitem{BB}
K. Bampouras and O. F. Brevig,
\emph{Norm attaining vectors and Hilbert points},
arXiv:2310.20346 (2023; revised 2024).

\bibitem{Brevig}
O. F. Brevig,
\emph{Minimal norm Hankel operators},
Proc. Amer. Math. Soc. \textbf{150} (2022), 4383--4391;
arXiv:2107.01680.

\bibitem{OS}
J. Ortega-Cerd\`a and K. Seip,
\emph{A lower bound in Nehari's theorem on the polydisc},
J. Anal. Math. \textbf{118} (2012), 339--342.
\end{thebibliography}

\end{document}
